\documentclass[10pt,a4paper]{amsart}
\usepackage[margin=19mm]{geometry}
\usepackage{amsmath,amssymb,mathtools}
\usepackage[scr=rsfs,cal=euler]{mathalfa}
\usepackage{xfrac}
\usepackage[unicode,hidelinks,bookmarks=false]{hyperref}
\newcommand{\ie}{i.e.\ }
\newcommand{\cf}{cf.\ }
\newcommand{\resp}{respectively\ }
\newcommand{\Subsection}{\subsection}
\providecommand{\nobreakdash}{\nobreak\hbox{-}}
\setlength{\emergencystretch}{2em}
\multlinegap0pt
\usepackage{mathtools}
\usepackage{enumerate}
\usepackage{amssymb}
\newtheorem{lemma}{Lemma}
\newcommand{\op}{\operatorname}
\title{On Algebraic Integrability\\ of Vector Fields \\ with Rational Function Coefficients\\ that Separate Variables}
\author{Przemys{\l}aw Grabowski}
\date{Source: arXiv:2609.00974v1 (CC BY 4.0)}
\begin{document}
\maketitle

\noindent\emph{Source and licence.} On Algebraic Integrability\\ of Vector Fields \\ with Rational Function Coefficients\\ that Separate Variables. Version arXiv:2609.00974v1 (CC BY 4.0), distributed by arXiv under the Creative Commons Attribution 4.0 International licence, http://creativecommons.org/licenses/by/4.0/, as recorded in the arXiv licence metadata for this exact version and shown on the abstract page https://arxiv.org/abs/2609.00974v1. Full licence notice: \texttt{licenses/arxiv-2609.00974-CC-BY-4.0.txt}.\par\smallskip

\noindent\textit{Editorial scope.} Counted unit: the article's lemma 1,2 (source0.tex lines 941--950). The article's complete original proof of it is delivered with it (source0.tex lines 952--976, one \texttt{proof} environment of 25 lines). No mathematics is written, repaired or re-derived: the statement and the proof are the article's own lines, in the article's own notation, and no retained line is substituted or re-flowed.  Retained uncounted as the source-local context the proof needs: lines 870--879.  Nothing was omitted from any retained span, so the package carries no omission record.  No retained line contains a citation, so the wrapper carries no bibliography and no bibliography entry.  No retained line contains a figure environment, an image-inclusion command or drawing code, so no image is delivered and the package declares no asset dependency.  Editorial formatting only, in addition: an AMS \texttt{amsart} preamble that restates the article's own declarations and macros used by the retained lines, namely the environments \texttt{lemma} on separate counters, together with the standard packages those lines need.  The retained environments print in this document as Lemma 1 in order of appearance, whereas the article prints the same items with the numbering of its own full document.  The complete native source of the article as distributed in the arXiv e-print for this exact version is bundled unchanged as \texttt{sources/209071/source0.tex}.
\begin{lemma}[Lagrange Approximation]\label{lem: Lagrange Approximation}
    Let $\alpha_1,\ldots,\alpha_d$ be distinct elements of a field $k$. Let $f(x)=c\prod_{i=1}^d(x-\alpha_i)$, where $c\in k$, $c\ne 0$. Let $g\in k[x]$ be a polynomial of degree $<d$, then
    \[
    g(x)=\sum_{i=1}^d g(\alpha_i)\frac{f(x)}{(x-\alpha_i)f'(\alpha_i)}.
    \]
    In particular, we have that 
    $
    [x^{d-1}]g(x)=c\sum_{i=1}^d \frac{g(\alpha_i)}{f'(\alpha_i)}
    $ holds. 
\end{lemma}

\begin{lemma}[Low Orders of Residue Subfields imply Explicit Numerators]\label{lem: explicit h for residues field of orders 1,2}
    Let $k$ be a field. Let $f\in k[x]$ be a separable irreducible polynomial over $k$ of degree $d$ with roots $a_1,\ldots,a_d$.  Let $h\in k[x]$ be a polynomial of degree $<d$.
    Then the field $L\coloneqq k\left(\frac{h(a_1)}{f'(a_1)},\frac{h(a_2)}{f'(a_2)} ,\frac{h(a_3)}{f'(a_3)}, \ldots, \frac{h(a_d)}{f'(a_d))}\right)$
    satisfies $(L:k) \le 2$ if and only if:
    \begin{enumerate}
        \item $(L:k) =1$, and $h=cf'$ for some $c\in k$.
        \item $(L:k) = 2$, and $f=g\cdot \overline{g}$ factorizes over $L$ into two conjugate factors, 
        \\and there is $c\in L\setminus k$ such that $h=cg' \overline{g}+ \overline{c}g\overline{g}'$.
    \end{enumerate}
\end{lemma}

\begin{proof}
    Let $F$ be the splitting field of $f$, i.e., $F=k(a_1,\ldots,a_d)$. The extension $F/k$ is Galois, because every splitting field of a separable polynomial is Galois. Let $G$ denote $\op{Gal}(F/k)$.

    $G$ acts transitively on all $a_i$, and thus on all $\frac{h(a_i)}{f(a_i)}$. Indeed, if $\sigma\in G$ satisfies $\sigma(a_1)=a_i$, then $\sigma(\frac{h(a_1)}{f(a_1)})=\frac{h(a_i)}{f(a_i)}$, because $h,f$ are polynomials over $k$.

    We put $\gamma_i=\frac{h(a_1)}{f(a_1)}$.
    
    We assume $(L:k)\le 2$. Thus, the Galois orbit of $\gamma_1\in L$ is the set $A=\{\gamma_1,\ldots,\gamma_d\}$, but, since the order of this extension is $\le 2$, this set is of size $1$, or $2$.

    If $|A|=1$, then for all $i$ we have $\gamma_1=\gamma_i$, thus $\gamma_1\in k$. By Lagrange approximation, Lemma \ref{lem: Lagrange Approximation}, applied to $h$ and $f'$, we get
    \begin{align*}
     h(x)&=\sum_{i=1}^d h(a_i)\frac{f(x)}{(x-a_i)f'(a_i)}=\sum_{i=1}^d \gamma_i\frac{f(x)}{(x-a_i)}=\\
     &=\gamma_1\sum_{i=1}^d \frac{f(x)}{(x-a_i)}=\gamma_1\sum_{i=1}^d f'(a_i)\frac{f(x)}{(x-a_i)f'(a_i)}=\gamma_1f'(x).
    \end{align*}
    We can take $c=\gamma_1$.

    If $|A|=2$, then the set $\{a_1,\ldots,a_d\}$ splits into two sets $A_1=\{a_i: \gamma_i=\gamma_1\}$ and $A_2=\{a_i: \gamma_i\ne\gamma_1\}$. By transitivity, these sets are of equal size $d/2$, thus $2|d$ and $d\ge 2$. We put $d=2d'$. We denote the roots in $A_1$ by $\alpha_1,\ldots,\alpha_{d'}$ and the ones in $A_2$ by $\beta_1,\ldots,\beta_{d'}$. We assume, without any loss of generality, that $a_2\in A_2$. So, we will write $\gamma_2$ for the second value.
    
    We put $g(x)=\prod_{i=1}^{d'}(x-\alpha_i)$ and $\overline{g}(x)=\prod_{i=1}^{d'}(x-\beta_i)$, so $f(x)=g(x)\overline{g}(x)$, if $f$ is monic. We observe that the subgroup of $G$ fixing $L$ fixes $g$ and $\overline{g}$, so they are polynomials over $L$ that are conjugate to each other with respect to $G$.

    The polynomial $H(x)=\gamma_1 g'(x)\overline{g}(x)+\gamma_2 g(x)\overline{g}'(x)-h(x)$ is by the definitions of degree $<d$ and zero at all roots of $f$, thus $H=0$, i.e., we can put $c=\gamma_1$ and we are done.

    A simple direct check proves the other implication. 
    This finishes the proof.
\end{proof}

\end{document}
